\fontsize{14}{15}\selectfont
En este trabajo de investigación nos enfocamos en abordar la problemática de la predicción de fallas de software, un desafío persistente en el campo de la ingeniería de software. La motivación detrás de esta investigación radica en la necesidad de mejorar la precisión en la identificación de errores en etapas tempranas del desarrollo de software, lo que a su vez contribuye a la optimización del uso de recursos y tiempo durante las fases de prueba y mantenimiento.

Para lograr este objetivo, se implementó un enfoque innovador que combina técnicas de selección de características, mediante análisis de relevancia y control de redundancia, con estrategias de balanceo de clases, utilizando el conjunto de datos BugHunter. Este enfoque fue diseñado para abordar desafíos específicos asociados con los conjuntos de datos de fallas, como la alta dimensionalidad, el desequilibrio de clases y la presencia de valores nulos.

Para responder a las preguntas de investigación planteadas, se realizaron experimentos exhaustivos que implicaron la aplicación de técnicas de selección de características y balanceo de clases en el conjunto de datos BugHunter, seguidos de la evaluación de los modelos de clasificación resultantes. La validación de la hipótesis se basó en la comparación de los resultados obtenidos con una línea base definida por el uso del conjunto de datos BugHunter sin procesar.

Los resultados experimentales confirmaron la efectividad de nuestro enfoque, demostrando mejoras significativas en la precisión de la predicción de fallas de software sin comprometer la integridad de los resultados. A través de un análisis detallado, se identificaron las métricas más relevantes para la predicción de errores y se demostró que ciertas métricas calculadas en BugHunter no influyen significativamente en la efectividad de la predicción.

La principal contribución de este trabajo radica en la demostración empírica de que la selección adecuada de características y el equilibrio de clases pueden tener un impacto positivo en la precisión de los modelos de predicción de fallas de software. Sin embargo, se identificaron limitaciones relacionadas con la generalización de los resultados a otros conjuntos de datos y contextos de software, así como la necesidad de investigar más a fondo el impacto de diferentes técnicas de balanceo de clases.

Basado en las limitaciones identificadas, el trabajo futuro podría explorar la aplicación de nuestro enfoque a una variedad más amplia de conjuntos de datos y contextos de software para validar su generalización. Además, se recomienda investigar en profundidad el impacto de otras técnicas de balanceo de clases y su interacción con diferentes tipos de modelos de clasificación, con el objetivo de refinar aún más la metodología propuesta para la predicción de fallas de software. Finalmente, una línea de investigación futura interesante sería la integración de estos datos en modelos de lenguaje de gran escala (LLM), con el fin de asistir en la detección automática de errores en fragmentos de código y potenciar la capacidad de análisis y predicción en entornos de desarrollo reales.



